\chapter*{Introduction Générale}
\addcontentsline{toc}{chapter}{Introduction Générale}

Face à l'avancée rapide des technologies digitales, telles que l'intelligence artificielle et l'Internet des objets, on assiste à l'émergence de nouvelles solutions intelligentes destinées à accompagner l'utilisateur dans ses tâches de tous les jours. Dans les contextes académiques et professionnels, les utilisateurs sont de plus en plus amenés à accomplir des activités qui requièrent de longues périodes de concentration, qu'il s'agisse d'étudier, de coder, d'analyser ou de travailler en général.

Toutefois, ces longues sessions de concentration intense peuvent causer divers problèmes tels que l'épuisement, les postures inappropriées, la diminution graduelle de la concentration et de la productivité ou encore les interruptions. Ces phénomènes se manifestent généralement graduellement et ne sont pas systématiquement décelés immédiatement par l'usager.

Malgré l'existence de certaines solutions dédiées à la productivité, celles-ci restent limitées car elles se concentrent principalement sur la gestion du temps ou le blocage des distractions numériques, sans analyser le comportement réel de l'utilisateur ni son état physique et cognitif. D'où la nécessité de concevoir un système intelligent capable d'observer et d'interpréter plusieurs indicateurs comportementaux afin d'évaluer de manière fiable le niveau de concentration.

Le projet \textbf{Smart Focus} a pour objectif de concevoir et développer un système intelligent capable d'observer, analyser et accompagner l'utilisateur durant ses sessions de travail ou d'étude. Il repose sur la combinaison de la vision par ordinateur, de l'intelligence artificielle et des technologies IoT afin de détecter des indicateurs tels que la posture, les signes de fatigue et les situations de distraction.

Le système vise à fournir une évaluation continue et fiable du niveau de concentration, accompagnée d'un feedback en temps réel, dans le but d'améliorer la productivité, la qualité du travail et le bien-être de l'utilisateur.

Pour la réalisation de ce projet, une méthodologie agile basée sur Scrum a été adoptée. Cette approche permet d'organiser le travail en sprints et d'assurer une progression incrémentale du système tout en facilitant l'adaptation aux besoins évolutifs.

Ce mémoire est organisé en plusieurs chapitres complémentaires, chacun abordant une étape spécifique du projet :

\begin{itemize}
    \item \textbf{Chapitre 1 : Cadre général du projet} \\
    Ce chapitre introduit le contexte du travail, précise la problématique étudiée et présente les besoins ainsi que l'étude de l'existant ayant orienté la conception du système.

    \item \textbf{Chapitre 2 : Sprint 1 – Authentification et interfaces de base} \\
    Ce chapitre présente la conception et la réalisation du module d’authentification, couvrant l’inscription, la connexion et la gestion du profil utilisateur côté backend et application mobile.
    Il décrit également la mise en place des interfaces principales de l’application — tableau de bord, statistiques, sommeil, paramètres — ainsi que la navigation entre les écrans.

    \item \textbf{Chapitre 3 : Sprint 2 – Surveillance intelligente de la concentration} \\
    Ce chapitre est consacré à l’implémentation du module de vision par ordinateur.
    Il détaille les techniques utilisées pour analyser l’état de l’utilisateur en temps réel à partir du flux vidéo, notamment la détection de la fatigue, des distractions et de la posture, ainsi que le calcul d’un score global de concentration transmis au backend.

    \item \textbf{Chapitre 4 : Sprint 3 – Planning intelligent et assistant d’apprentissage} \\
    Ce chapitre décrit le développement du système de planning intelligent basé sur l’intelligence artificielle pour la génération automatique des sessions d’étude.
    Il présente également l’intégration du chatbot RAG permettant l’interrogation de documents PDF, la génération automatique de quiz et de flashcards avec révision espacée (SM-2), ainsi que la connexion des interfaces mobiles aux données dynamiques du backend.

    \item \textbf{Chapitre 5 : Sprint 4 – Système IoT et alertes intelligentes} \\
    Ce chapitre présente l’intégration du dispositif matériel de monitoring permettant la capture continue des données utilisateur via caméra et capteurs.
    Il décrit également la mise en place du suivi en temps réel dans l’application mobile ainsi que du système d’alertes intelligentes basé sur le niveau de concentration et les données de sommeil, afin d’adapter le comportement du système et le planning de travail.
\end{itemize}

Tout au long de ce travail, nous mettons en évidence les choix réalisés, les contraintes rencontrées ainsi que les solutions adoptées pour mener à bien ce projet.

\clearpage

%%%%%%%%%%%%%%%%%%%%%%%%%%%%%%